% =====================================================================
\section{ODMR line shapes}
\label{app:odmr}
% =====================================================================

This appendix derives the line shapes summarized in Sec.~\ref{sec:model}: the CW line with the laser treated as two
rates rather than a reset, and the pulsed line with its inhomogeneous broadening and readout.

\paragraph{CW line with optical pumping.} Eq.~\eqref{eq:cw-lorentzian} treats the laser only as a reset, which is too
crude when the resets are frequent. The laser acts on the spin through two rates, both set by how often the center
goes through an optical cycle: it pumps the spin back toward $m_s=0$, at a rate $\Gamma_p$, and each cycle destroys
the phase between $m_s=0$ and $m_s=-1$, at a rate $\Gamma_c$. The microwave enters only through its angular Rabi
frequency $\Omega_R=2\pi\Omega$. The competition between the two is clearest in terms of populations. A drive whose
coherence is destroyed at the rate $\Gamma_c$ no longer produces Rabi oscillations; it moves population between
$m_s=0$ and $m_s=-1$, in both directions alike, at the rate
\begin{equation}
  W(\Delta)=\frac{\Omega_R^2\,\Gamma_c}{2\big[\Gamma_c^2+(2\pi\Delta)^2\big]}
  =\frac{\Omega_R^2}{2\Gamma_c}\,\frac{\Gamma_c^2}{\Gamma_c^2+(2\pi\Delta)^2},
  \label{eq:transfer-rate}
\end{equation}
a Lorentzian of peak $\Omega_R^2/2\Gamma_c$ and half width $\Gamma_c/2\pi$ in Hz. The laser returns population from
$m_s=-1$ at the rate $\Gamma_p$, so with $P_0=1-P_{-1}$ the population of $m_s=-1$ obeys
\begin{equation*}
  \dot P_{-1}=W\,(P_0-P_{-1})-\Gamma_pP_{-1}=W-(2W+\Gamma_p)\,P_{-1}.
\end{equation*}
In steady state, inserting Eq.~\eqref{eq:transfer-rate} and then ordering the terms of the denominator,
\begin{equation}
  \begin{aligned}
  P_{-1}^{\rm ss}(\Delta)&=\frac{W}{2W+\Gamma_p}
  =\frac12\,\frac{\Omega_R^2\,\Gamma_c}{\Omega_R^2\,\Gamma_c+\Gamma_p\big[\Gamma_c^2+(2\pi\Delta)^2\big]}\\
  &=\frac12\,\frac{\Omega_R^2\,\Gamma_c/\Gamma_p}{(2\pi\Delta)^2+\Gamma_c^2+\Omega_R^2\,\Gamma_c/\Gamma_p},
  \end{aligned}
  \label{eq:bloch-ss}
\end{equation}
the same as the steady state of the full Bloch equations \cite{dreau2011}. It reduces to
Eq.~\eqref{eq:cw-lorentzian} when the two rates are equal and small against $\Omega_R$. The last form is a
Lorentzian in $2\pi\Delta$ whose squared half width is the constant in its denominator, which gives
Eq.~\eqref{eq:linewidth}.

The first term of Eq.~\eqref{eq:linewidth} is the width the laser alone imposes, through the loss of coherence. The
second is power broadening by the microwave: a strong drive saturates the transfer, $W\gg\Gamma_p$, already on the
tails of the line, and the line widens until $W$ there has fallen to about the pumping rate. The depth of the line,
Eq.~\eqref{eq:bloch-ss} at $\Delta=0$, is
\begin{equation*}
  P_{-1}^{\rm ss}(0)=\frac12\,\frac{\Omega_R^2}{\Omega_R^2+\Gamma_p\Gamma_c}:
\end{equation*}
it grows with the microwave power, toward $\tfrac12$, and falls with the laser power, which repolarizes the spin
before the drive can flip it.

\paragraph{Optical saturation.} The two laser rates grow with the laser power, but not without limit. The optical
transition is itself a two-level system: the laser excites the center at a rate proportional to its intensity
$I_{\rm opt}$, and the excited state decays in its lifetime $\tau_e$. In steady state the excited population is
$s/(1+s)$, with the saturation parameter $s=I_{\rm opt}/I_{\rm sat}$ the ratio of the excitation rate to the decay
rate $1/\tau_e$. At low power it grows linearly with $s$; at high power it levels off, because a center cannot be
excited again before it has decayed. The rate of optical cycles is this population divided by $\tau_e$, and both
laser rates follow it, $\Gamma_{p,c}=\Gamma_{p,c}^\infty\,s/(1+s)$ \cite{dreau2011}, with $\Gamma_{p,c}^\infty$ their
values at full saturation. Raising the laser power therefore both widens and flattens the line, up to saturation.
% Dreau et al. (2011) Eq. 11 quotes the same width as Gamma_c^inf/(2 pi) sqrt((s/(1+s))^2 + Omega_R^2/(Gamma_p^inf
% Gamma_c^inf)) and calls it the FWHM; with Gamma_{p,c} = Gamma^inf s/(1+s) this is Eq. (linewidth) up to their
% factor-2 convention for the rates. Check before quoting numbers.

\paragraph{Pulsed line.} The $\pi$-pulse length is chosen so that on resonance, where the spin oscillates at
$\Omega$, it is fully inverted: $\pi\Omega t_\pi=\pi/2$. Off resonance it oscillates at $\Omega_{\rm gen}$ instead, so
the argument of the sine in Eq.~\eqref{eq:rabi} at $t=t_\pi$ is $\pi\Omega_{\rm gen} t_\pi=(\pi/2)\,\Omega_{\rm gen}/\Omega$,
which gives Eq.~\eqref{eq:pulsed}. On resonance $\Omega_{\rm gen}=\Omega$ and the inversion is complete. For
$|\Delta|\gg\Omega$, $\Omega_{\rm gen}\approx|\Delta|$ and
$P_{-1}^{\pi}\approx\sin^2(\pi\Delta t_\pi)/(2\Delta t_\pi)^2$, the power spectrum of a rectangular pulse of length
$t_\pi$: far from resonance the spin responds to the frequency content of the pulse. The line vanishes wherever
$\Omega_{\rm gen}=2\Omega,4\Omega,\dots$, i.e.\ at $|\Delta|=\sqrt3\,\Omega,\sqrt{15}\,\Omega,\dots$: there the spin
completes a whole number of full Rabi cycles during the pulse and returns to $m_s=0$. Between the zeros lie side
lobes that decay as $\Omega^2/\Omega_{\rm gen}^2$; the first, near $\Omega_{\rm gen}=3\Omega$ at $|\Delta|\approx2.7\,\Omega$,
is 12\% of the depth.

\paragraph{Inhomogeneous broadening.} The transition frequency of the spin is not sharp: slowly fluctuating local
fields, mostly from the bath of $^{13}$C nuclear spins, shift it by a random offset $\xi$ that is fixed during one
pulse but changes from one repetition of the sequence to the next. A measurement averages many repetitions, each
driven at the detuning $\Delta-\xi$ instead of $\Delta$, so the measured line is Eq.~\eqref{eq:pulsed} averaged over
the distribution of $\xi$, i.e.\ convolved with it. For a Gaussian distribution of standard deviation $\sigma_\xi$,
\begin{equation}
  \bar P_{-1}^{\pi}(\Delta)=\int P_{-1}^{\pi}(\Delta-\xi)\,\frac{e^{-\xi^2/2\sigma_\xi^2}}{\sqrt{2\pi}\,\sigma_\xi}\,d\xi .
  \label{eq:pulsed-t2star}
\end{equation}
The width $\sigma_\xi$ is what $T_2^*$ measures. A free spin precessing at the offset $\xi$ picks up the phase
$2\pi\xi t$, and the average of $e^{2\pi i\xi t}$ over the Gaussian is the coherence envelope
$e^{-2\pi^2\sigma_\xi^2t^2}$. $T_2^*$ is defined by writing this envelope as $e^{-(t/T_2^*)^2}$, so
$\sigma_\xi=1/(\sqrt2\,\pi T_2^*)$. For $t_\pi\ll T_2^*$ the spread is small against $\Omega$ and the line is that of
Eq.~\eqref{eq:pulsed}. For $t_\pi\gtrsim T_2^*$ the Gaussian dominates: the zeros and side lobes are washed out, the
line stops narrowing at a width $\sim1/T_2^*$, and the depth falls because the pulse no longer inverts every
frequency class \cite{dreau2011}.

\paragraph{Readout.} While the microwave pulse is on, the laser is off and the spin relaxes only through $T_1$, which
is much longer than $t_\pi$, and dephases through the spread already included in Eq.~\eqref{eq:pulsed-t2star}. The
readout is different: the laser that reads the spin out also repolarizes it, so the PL difference between $m_s=0$
and $m_s=-1$ lasts only for the repolarization time $\tau_p$. Taking the repolarization as a single exponential, a
spin flipped with probability $P_{-1}$ emits, a time $t$ into the readout pulse, PL at the rate
$\Phi_0\,[1-c\,P_{-1}\,e^{-t/\tau_p}]$, where $\Phi_0$ is the rate of $m_s=0$ and $c$ the fractional PL drop of
$m_s=-1$ at the start of the readout. The recorded PL is this rate integrated over a readout window of length $t_r$,
normalized to the same integral with no flip,
\begin{equation}
  \begin{aligned}
  \mu&=1-\frac{1}{t_r}\int_0^{t_r}c\,P_{-1}\,e^{-t/\tau_p}\,dt\\
  &=1-c\,\frac{\tau_p}{t_r}\big(1-e^{-t_r/\tau_p}\big)\,\bar P_{-1}^{\pi}(\Delta) ,
  \end{aligned}
  \label{eq:pulsed-readout}
\end{equation}
with $P_{-1}=\bar P_{-1}^{\pi}$ of Eq.~\eqref{eq:pulsed-t2star}. The integration leaves the line shape unchanged and
sets only its depth: the contrast is $C=c\,(\tau_p/t_r)(1-e^{-t_r/\tau_p})$. A longer window collects more photons,
$N_{\rm ph}\propto t_r$, but dilutes the contrast, so the single-shot signal-to-noise ratio
$C\sqrt{N_{\rm ph}}\propto(1-e^{-t_r/\tau_p})/\sqrt{t_r/\tau_p}$ is largest at $t_r\approx1.26\,\tau_p$.
